\section{Related Work}
\label{sec:related}

\paragraph{Popularity, social learning and herding.}
When quality is hard to observe, adopters rationally follow earlier adopters
\citep{banerjee1992simple,bikhchandani1992theory}, and social influence makes success both more unequal and less
predictable from quality \citep{salganik2006experimental}. Cumulative advantage in science has the same structure
\citep{merton1968matthew}. Field experiments show that displaying popularity information shifts choices
\citep{cai2009observational,tucker2011popularity}. In observational data, the central difficulty is that popularity
is itself a signal of quality. We address it in three ways: by bounding quality with product-specific fixed effects,
the largest share any unobserved fixed attribute could take; by measuring how much of that bound external benchmark
scores account for; and by comparing identical copies of the same weights, where quality is held fixed by
construction.

\paragraph{Demand estimation with very large choice sets.}
Conditional logit \citep{mcfadden1974conditional,train2009discrete} remains consistent when the likelihood is
evaluated on a sample of the alternatives, provided the sampling probabilities enter as an offset
\citep{mcfadden1978modelling}; related results cover semiparametric models \citep{fox2007semiparametric} and
multivariate extreme-value models \citep{guevara2013sampling}. These results allow any positive sampling
probabilities. Our setting, with hundreds of thousands of alternatives whose demand is concentrated on a few dozen,
makes the choice of design consequential. We sample with probability proportional to degree plus one, so that the
incumbents developers actually weigh are present in every choice set, and we show that a uniform design changes the
estimated weights and does not only widen the intervals.

\paragraph{Network formation as discrete choice.}
Preferential attachment explains heavy-tailed degree distributions by new nodes linking to existing nodes in
proportion to their degree \citep{price1976general,barabasi1999emergence}. Fitness models add a node-specific
attractiveness that lets late entrants overtake incumbents \citep{bianconi2001competition,wang2013quantifying}, and
empirical work measures the attachment kernel from the timing of new links
\citep{newman2001clustering,jeong2003measuring}. In economics, \citet{jackson2007meeting} show how a mix of random
and network-based search produces such degree distributions, and econometric models of network formation estimate
link preferences from observed networks \citep{graham2017econometric,depaula2020econometric}.
\citet{overgoor2019choosing} show that the physics mechanisms are conditional-logit models in which each new edge is a
choice among the existing nodes, so that degree, fitness, homophily and triadic closure can be estimated jointly and
compared by likelihood; importance sampling extends this to large graphs \citep{overgoor2020scaling}. We build on this
framework. We compare mechanisms with a Shapley decomposition of the likelihood gain instead of a sequence of nested
models, and we use the estimated utility to grow counterfactual networks.

\paragraph{AI model ecosystems.}
Foundation models are adapted rather than trained from scratch \citep{bommasani2021opportunities}, which makes the
choice of a base model a consequential and recurring decision. Descriptive studies of the Hugging Face Hub document
its growth, its family trees and the concentration of downstream use on a few model families
\citep{osborne2024ai,laufer2025anatomy}. In machine learning, base-model choice is usually framed as a prediction
problem: transferability estimators such as LEEP and LogME rank pre-trained models by how well they are expected to
fine-tune on a target task \citep{nguyen2020leep,you2021logme}. Our question is complementary and behavioral. We ask
not which model a developer \emph{should} choose but which signals the choices developers actually make respond to,
and what kind of market those choices grow.
